% ========================================================
% BÀI TẬP TOÁN 5 - KẾT NỐI TRI THỨC VỚI CUỘC SỐNG
% Bài 26: Hình thang. Diện tích hình thang
% Định dạng: BTTeX Part III - Trắc nghiệm điền số ngắn
% ========================================================

% Câu 1
\begin{ex}
  Một hình thang có đáy lớn $15\text{ cm}$, đáy bé $9\text{ cm}$ và chiều cao $6\text{ cm}$. Diện tích hình thang đó là bao nhiêu xăng-ti-mét vuông?
  \shortans{72}
  \loigiai{
    $(15 + 9) \times 6 : 2 = 24 \times 3 = 72\text{ cm}^2$.
  }
\end{ex}

% Câu 2
\begin{ex}
  Một hình thang có đáy lớn $20\text{ m}$, đáy bé $12\text{ m}$ và chiều cao $10\text{ m}$. Diện tích hình thang đó là bao nhiêu mét vuông?
  \shortans{160}
  \loigiai{
    $(20 + 12) \times 10 : 2 = 160\text{ m}^2$.
  }
\end{ex}

% Câu 3
\begin{ex}
  Hình thang có tổng hai đáy là $34\text{ cm}$ và chiều cao là $8\text{ cm}$. Diện tích hình thang đó là bao nhiêu xăng-ti-mét vuông?
  \shortans{136}
  \loigiai{
    $(34 \times 8) : 2 = 136\text{ cm}^2$.
  }
\end{ex}

% Câu 4
\begin{ex}
  Hình thang có diện tích $90\text{ m}^2$, chiều cao $6\text{ m}$. Tổng độ dài hai đáy là bao nhiêu mét?
  \shortans{30}
  \loigiai{
    Tổng hai đáy là: $(90 \times 2) : 6 = 30\text{ m}$.
  }
\end{ex}

% Câu 5
\begin{ex}
  Hình thang có diện tích $120\text{ cm}^2$, tổng hai đáy là $24\text{ cm}$. Chiều cao của hình thang là bao nhiêu xăng-ti-mét?
  \shortans{10}
  \loigiai{
    Chiều cao là: $(120 \times 2) : 24 = 10\text{ cm}$.
  }
\end{ex}

% Câu 6
\begin{ex}
  Một hình thang vuông có chiều cao bằng cạnh bên vuông góc là $7\text{ cm}$. Biết tổng hai đáy là $20\text{ cm}$. Diện tích hình thang vuông là bao nhiêu xăng-ti-mét vuông?
  \shortans{70}
  \loigiai{
    $(20 \times 7) : 2 = 70\text{ cm}^2$.
  }
\end{ex}

% Câu 7
\begin{ex}
  Một thửa ruộng hình thang có đáy lớn $50\text{ m}$, đáy bé $30\text{ m}$, chiều cao $25\text{ m}$. Diện tích thửa ruộng là bao nhiêu mét vuông?
  \shortans{1000}
  \loigiai{
    $(50 + 30) \times 25 : 2 = 80 \times 25 : 2 = 1\,000\text{ m}^2$.
  }
\end{ex}

% Câu 8
\begin{ex}
  Ở thửa ruộng trên, trung bình mỗi $100\text{ m}^2$ thu hoạch được $60\text{ kg}$ thóc. Cả thửa ruộng thu hoạch được bao nhiêu tạ thóc?
  \shortans{6}
  \loigiai{
    Số thóc là: $(1\,000 : 100) \times 60 = 600\text{ kg} = 6\text{ tạ}$.
  }
\end{ex}

% Câu 9
\begin{ex}
  Một hình thang có đáy bé $10\text{ cm}$, đáy lớn gấp đôi đáy bé, chiều cao bằng $8\text{ cm}$. Diện tích hình thang là bao nhiêu xăng-ti-mét vuông?
  \shortans{120}
  \loigiai{
    Đáy lớn là $20\text{ cm}$. Diện tích là $(10 + 20) \times 8 : 2 = 120\text{ cm}^2$.
  }
\end{ex}

% Câu 10
\begin{ex}
  Một hình thang có đáy lớn $18\text{ cm}$, đáy bé kém đáy lớn $6\text{ cm}$, chiều cao bằng trung bình cộng hai đáy. Chiều cao hình thang là bao nhiêu xăng-ti-mét?
  \shortans{15}
  \loigiai{
    Đáy bé là $12\text{ cm}$. Chiều cao là $(18 + 12) : 2 = 15\text{ cm}$.
  }
\end{ex}

% Câu 11
\begin{ex}
  Ở câu trên, diện tích hình thang đó là bao nhiêu xăng-ti-mét vuông?
  \shortans{225}
  \loigiai{
    Diện tích là $15 \times 15 = 225\text{ cm}^2$.
  }
\end{ex}

% Câu 12
\begin{ex}
  Một hình thang có diện tích $84\text{ cm}^2$, chiều cao $7\text{ cm}$, đáy lớn hơn đáy bé $4\text{ cm}$. Độ dài đáy lớn là bao nhiêu xăng-ti-mét?
  \shortans{14}
  \loigiai{
    Tổng hai đáy là $(84 \times 2) : 7 = 24\text{ cm}$. Đáy lớn là $(24 + 4) : 2 = 14\text{ cm}$.
  }
\end{ex}

% Câu 13
\begin{ex}
  Ở câu trên, độ dài đáy bé là bao nhiêu xăng-ti-mét?
  \shortans{10}
  \loigiai{
    Đáy bé là $24 - 14 = 10\text{ cm}$.
  }
\end{ex}

% Câu 14
\begin{ex}
  Một mảnh đất hình thang có đáy lớn $40\text{ m}$, đáy bé $20\text{ m}$. Nếu mở rộng đáy lớn thêm $5\text{ m}$ thì diện tích tăng thêm $40\text{ m}^2$. Chiều cao của mảnh đất là bao nhiêu mét?
  \shortans{16}
  \loigiai{
    Chiều cao là: $(40 \times 2) : 5 = 16\text{ m}$.
  }
\end{ex}

% Câu 15
\begin{ex}
  Ở câu trên, diện tích mảnh đất ban đầu là bao nhiêu mét vuông?
  \shortans{480}
  \loigiai{
    Diện tích ban đầu: $(40 + 20) \times 16 : 2 = 480\text{ m}^2$.
  }
\end{ex}

% Câu 16
\begin{ex}
  Đáy lớn của hình thang là $24\text{ cm}$, đáy bé bằng $\frac{2}{3}$ đáy lớn. Chiều cao bằng $10\text{ cm}$. Diện tích hình thang là bao nhiêu xăng-ti-mét vuông?
  \shortans{200}
  \loigiai{
    Đáy bé là $16\text{ cm}$. Diện tích là $(24 + 16) \times 10 : 2 = 200\text{ cm}^2$.
  }
\end{ex}

% Câu 17
\begin{ex}
  Một hình thang có đáy lớn $16\text{ cm}$, đáy bé $8\text{ cm}$. Chiều cao bằng nửa đáy lớn. Diện tích hình thang là bao nhiêu xăng-ti-mét vuông?
  \shortans{96}
  \loigiai{
    Chiều cao là $8\text{ cm}$. Diện tích là $(16 + 8) \times 8 : 2 = 96\text{ cm}^2$.
  }
\end{ex}

% Câu 18
\begin{ex}
  Hình thang cân có hai cạnh bên bằng nhau và mỗi cạnh dài $10\text{ cm}$. Hai đáy dài lần lượt là $12\text{ cm}$ và $18\text{ cm}$. Chu vi hình thang cân là bao nhiêu xăng-ti-mét?
  \shortans{50}
  \loigiai{
    Chu vi là $12 + 18 + 10 + 10 = 50\text{ cm}$.
  }
\end{ex}

% Câu 19
\begin{ex}
  Một khúc đê có mặt cắt là hình thang với đáy lớn $12\text{ m}$, đáy bé $4\text{ m}$ và chiều cao $3\text{ m}$. Diện tích mặt cắt khúc đê là bao nhiêu mét vuông?
  \shortans{24}
  \loigiai{
    $(12 + 4) \times 3 : 2 = 24\text{ m}^2$.
  }
\end{ex}

% Câu 20
\begin{ex}
  Nếu gấp cả hai đáy và chiều cao của hình thang lên $2$ lần thì diện tích hình thang gấp lên mấy lần?
  \shortans{4}
  \loigiai{
    Diện tích tăng $2 \times 2 = 4$ lần.
  }
\end{ex}

% Câu 21
\begin{ex}
  Một biển báo giao thông hình thang có diện tích $36\text{ dm}^2$, chiều cao $6\text{ dm}$. Tổng hai đáy của biển báo là bao nhiêu đề-xi-mét?
  \shortans{12}
  \loigiai{
    Tổng hai đáy là $(36 \times 2) : 6 = 12\text{ dm}$.
  }
\end{ex}

% Câu 22
\begin{ex}
  Một thửa ruộng hình thang có diện tích $600\text{ m}^2$, chiều cao $20\text{ m}$. Đáy lớn gấp đôi đáy bé. Đáy bé dài bao nhiêu mét?
  \shortans{20}
  \loigiai{
    Tổng hai đáy là $(600 \times 2) : 20 = 60\text{ m}$. Đáy bé là $60 : (1 + 2) = 20\text{ m}$.
  }
\end{ex}

% Câu 23
\begin{ex}
  Ở câu trên, đáy lớn dài bao nhiêu mét?
  \shortans{40}
  \loigiai{
    Đáy lớn là $20 \times 2 = 40\text{ m}$.
  }
\end{ex}

% Câu 24
\begin{ex}
  Một hình thang có đáy lớn $15\text{ cm}$, đáy bé $7\text{ cm}$. Cắt hình thang theo một đường thẳng song song với hai đáy chia hình thang thành hai hình thang nhỏ. Tổng diện tích hai hình thang nhỏ so với diện tích hình thang ban đầu bằng mấy lần?
  \shortans{1}
  \loigiai{
    Tổng diện tích hai hình nhỏ bằng đúng diện tích ban đầu (1 lần).
  }
\end{ex}

% Câu 25
\begin{ex}
  Cho hình thang $ABCD$ ($AB // CD$). Nối hai đường chéo $AC$ và $BD$ cắt nhau tại $O$. Biết diện tích tam giác $AOD$ là $15\text{ cm}^2$. Diện tích tam giác $BOC$ là bao nhiêu xăng-ti-mét vuông?
  \shortans{15}
  \loigiai{
    Trong hình thang, hai tam giác ở hai cạnh bên có diện tích bằng nhau: $S_{BOC} = S_{AOD} = 15\text{ cm}^2$.
  }
\end{ex}
